
\subsection{电势与能}

\begin{tabularx}{\columnwidth}{XX}
  电势定义&$\varphi=\dfrac{E_p}{q}$\\
  电势差定义&$U_{AB}=\varphi_A-\varphi_B$\\
  电场力做功&$W_{AB}=qU_{AB}$\\
  电势能变化&$\Delta E_p=-W_{\text{电}}$\\
  匀强电场电势差&$U=Ed$\\
  电场力做功&$W=qEd\cos\theta$\\
\end{tabularx}

\begin{formulanotebox}
$U=Ed$ 只适用于匀强电场，且 $d$ 是沿电场方向的投影距离。若位移与电场方向成角度，必须先取投影。

电势是位置属性，电势能还与电荷正负有关。正电荷在高电势处电势能大，负电荷相反。
\end{formulanotebox}

\begin{keypointbox}[title={\strut\textcolor{white}{\bfseries 重点知识点：电势与做功}}]
\textbf{做功路径无关}：静电力做功只与初末位置有关，与路径无关，因此静电场是保守场。\\

\textbf{能量判断}：电场力做正功，电势能减少；电场力做负功，电势能增加。若只有电场力做功，动能和电势能相互转化，总机械能守恒。\\

\textbf{等势面}：电场线与等势面垂直，沿等势面移动电荷时电场力不做功。
\end{keypointbox}

\begin{casetypebox}[title={\strut\textcolor{white}{\bfseries 常见题型：电势高低与能量变化}}]
\begin{itemize}[leftmargin=1.5em]
  \item 先根据电场线方向判断电势降低方向，再结合电荷正负判断电势能变化。
  \item 用 $W_{AB}=qU_{AB}$ 时，$q$ 要带符号代入。
  \item 带电粒子从静止经电势差 $U$ 加速时，常用 $qU=\frac12mv^2$。
  \item 匀强电场中若给出两点距离而不是投影距离，应先画出沿电场方向的有效距离。
\end{itemize}
\end{casetypebox}
